% 小行星（综述）
% license CCBYSA3
% type Wiki

本文根据 CC-BY-SA 协议转载翻译自维基百科\href{https://en.wikipedia.org/wiki/Asteroid}{相关文章}。

\begin{figure}[ht]
\centering
\includegraphics[width=8cm]{./figures/9fdf35f418301f87.png}
\caption{被探测访问过的小行星影像展示其差异：（上排）433 Eros 与 243 Ida 及其卫星 Dactyl，（下排）Ceres 与 101955 Bennu。尺寸不按比例绘制。} \label{fig_Astero_1}
\end{figure}
小行星是一类小行星体（minor planet），其尺寸大于流星体（meteoroid），但既不是行星，也不是已确认为彗星的天体。它们位于太阳系内部轨道运行，或与木星共享轨道（如特洛伊小行星）。小行星由岩石、金属或冰质物质构成，不具备大气层，并可大致分类为 C 型（碳质）、M 型（金属）、S 型（硅酸盐）。小行星的大小与形状差异极大，从直径不足一千米的松散碎石堆，到直径接近 1000 千米、被归类为矮行星的 Ceres。若天体在受到太阳辐射加热时产生彗发（尾），则被归类为彗星而非小行星，尽管近期观测显示这两类天体之间可能存在连续光谱。\(^\text{[1][2]}\)

在约一百万颗已知小行星中，\(^\text{[3]}\)数量最多的位于火星与木星轨道之间、距离太阳约 2 至 4 AU 的区域，称为主小行星带（main asteroid belt）。所有小行星总质量仅约为地球月球质量的 3\%。主带中大多数小行星沿轻微椭圆、稳定的轨道运行，方向与地球一致，绕太阳一周约需 3 至 6 年。\(^\text{[4]}\)

小行星最早通过地面观测方式进行研究，首次近距离观测来自伽利略号（Galileo）探测器。其后，美国 NASA 与日本 JAXA 相继发射了多次专门的小行星探测任务，并持续规划新的任务。NASA 的 NEAR Shoemaker 研究了 Eros，Dawn 探测了 Vesta 与 Ceres；JAXA 的隼鸟号（Hayabusa）与隼鸟 2 号分别研究并回收了 Itokawa 与 Ryugu 的样本；OSIRIS{-}REx 探测 Bennu，并于 2020 年采样，2023 年将样本送回地球。NASA 的 Lucy（2021 年发射）将研究十颗不同小行星，其中两颗来自主带、八颗为木星特洛伊小行星。Psyche 探测器于 2023 年 10 月发射，目标是金属小行星 Psyche。欧洲航天局的 Hera 于 2024 年 10 月发射，旨在研究 DART 碰撞后的结果。中国国家航天局的天问二号于 2025 年 5 月发射，\(^\text{[5]}\)其任务为探测共轨近地小行星 469219 Kamoʻoalewa 与活跃小行星 311P/PanSTARRS，并对 Kamoʻoalewa 的风化层（regolith）进行采样。\(^\text{[6]}\)

近地小行星若与地球发生撞击，可能造成灾难性影响，其中最著名的案例是 Chicxulub 撞击事件，被普遍认为导致了白垩纪–古近纪生物大灭绝。为应对这一潜在威胁，作为实验性防御措施，2022 年 9 月，双小行星偏转测试（Double Asteroid Redirection Test, DART）探测器通过撞击成功改变了无威胁小行星 Dimorphos 的轨道。
\subsection{术语}
\begin{figure}[ht]
\centering
\includegraphics[width=8cm]{./figures/795a52138582bb97.png}
\caption{一幅按相同尺度制作的组合图像，展示了 2012 年前以高分辨率成像的小行星。从大到小依次为：4 Vesta、21 Lutetia、253 Mathilde、243 Ida 及其卫星 Dactyl、433 Eros、951 Gaspra、2867 Šteins、25143 Itokawa。} \label{fig_Astero_2}
\end{figure}
\begin{figure}[ht]
\centering
\includegraphics[width=8cm]{./figures/1f1fc84d80338171.png}
\caption{按相同比例展示的 Vesta（左）、Ceres（中）与月球（右）} \label{fig_Astero_3}
\end{figure}
2006 年，国际天文学联合会（IAU）引入了当前更为推荐的宽泛术语“太阳系小天体”（small Solar System body），将其定义为太阳系中既不是行星、也不是矮行星、亦不是天然卫星的天体；其中包括小行星、彗星以及更近期发现的其他类型天体\(^\text{[7]}\)。根据 IAU 表述，“术语『小行星（minor planet）』仍可继续使用，但一般情况下将优先采用『太阳系小天体（Small Solar System Body）』”\(^\text{[8]}\)。

在历史上，第一颗被发现的小行星——Ceres——最初曾被视为一颗新行星\(^\text{[a]}\)。随后，天文学家又陆续发现了其他相似天体。在当时的观测设备条件下，它们呈现为类似恒星的光点，几乎看不到行星盘结构，但由于具有明显的视运动，又可与恒星区分开来。这促使天文学家威廉·赫歇尔（Sir William Herschel）提出了“asteroid”一词\(^\text{[b]}\)，该词源自希腊语 ἀστεροειδής（asteroeidēs），意为“类星”“似星状”，并由古希腊语 ἀστήρ（astēr，“星、行星”）派生。在 19 世纪后半叶初期，“asteroid”与“planet”（有时并未特别限定为“minor”）仍被交替使用\(^\text{[c]}\)。

传统上，围绕太阳运行的小型天体被划分为彗星、小行星或流星体，其中直径小于一米的天体被称为流星体。术语“asteroid”从未被正式定义\(^\text{[13]}\)，但在非正式使用中可指“围绕太阳运行、形状不规则、呈岩质组成且不符合 IAU 定义中行星或矮行星条件的天体”\(^\text{[14]}\)。小行星与彗星的主要区别在于，彗星因近表层冰在太阳辐射作用下发生升华而形成彗发（尾）。少数天体最初被归类为小行星，但后续观测显示存在彗状活动。相反，一些（或许全部）彗星最终会耗尽其表面挥发性冰，从而演化为类小行星天体。另一区分方式是轨道偏心率：彗星通常具有比大多数小行星更高的轨道偏心率；而高度偏心的小行星极可能为休眠或灭绝的彗星\(^\text{[15]}\)。

位于木星轨道之外的小行星体有时亦被称为“asteroids”，尤其是在大众化科普材料中\(^\text{[d]}\)。然而，将“asteroid”一词限定用于太阳系内部的小行星体的做法正在变得愈发普遍\(^\text{[17]}\)。因此，本文主要将讨论范围限定于传统意义上的小行星：即小行星带天体、木星特洛伊小行星以及近地天体。

自 1801 年 Ceres 被发现后的近两个世纪中，所有已知小行星在其轨道的大部分时间里都位于木星轨道以内，尽管极少数如 944 Hidalgo 会在轨道部分区段上延伸至更远处。自 1977 年 2060 Chiron 被发现起，天文学家开始发现永久位于木星轨道之外的小型天体，即后称为半人马（centaurs）。1992 年，又发现了 15760 Albion，这是首个已知位于海王星轨道之外的天体（除 Pluto 外）；随后人们观测到了大量类似天体，即现称为跨海王星天体（trans-Neptunian object）。更外层依次还有柯伊伯带天体、散射盘天体以及距离更为遥远的奥尔特云——其被假设为休眠彗星的主要储库。它们位于太阳系寒冷的外缘区域，冰质物质能够长期保持固态，因此彗状活动极为微弱；若半人马或跨海王星天体接近太阳，其挥发性冰质材料将发生升华，按传统分类方法将被归类为彗星。

之所以将柯伊伯带天体称为“objects”，部分原因是为了避免必须将其归为小行星或彗星\(^\text{[1]}\)。它们被认为在成分上大体类似于彗星，但部分天体可能与小行星更相近\(^\text{[18]}\)。绝大多数此类天体并不具有典型彗星所表现的高偏心率轨道，且目前发现者均比传统彗核更大。近年来的观测结果，如 Stardust 探测器采集的彗星尘埃样本分析，正不断模糊小行星与彗星之间的界限\(^\text{[1]}\)，越来越多研究表明二者之间可能形成连续谱，而非清晰的分类边界\(^\text{[2]}\)。

2006 年，IAU 为最大的小行星体引入了“矮行星”这一类别——即因自身引力而形成近似椭球形状的天体。迄今为止，仅小行星带中体积最大的天体 Ceres（直径约 975 km，约 606 mi）被划入此类\(^\text{[19][20]}\)。
\subsection{观测历史}
尽管小行星数量极为庞大，但其发现在人类天文学史上相对较晚，第一颗小行星——Ceres——直到 1801 年才被识别出来\(^\text{[21]}\)。在暗夜且观测条件良好的情况下，只有一颗小行星，即表面反照率较高的 4 Vesta，通常能够以肉眼看见。极少数情况下，当小型小行星非常接近地球时，也可能在短暂时间内被肉眼观测到\(^\text{[22]}\)。截至 2025 年 5 月，小行星中心（Minor Planet Center）已记录太阳系内外共计 1,460,356 颗小行星的数据，其中约有 826,864 颗具备足够观测信息而获得编号编号\(^\text{[23]}\)。
\subsubsection{Ceres 的发现}
\begin{figure}[ht]
\centering
\includegraphics[width=6cm]{./figures/63dcf87fd16a81c7.png}
\caption{《Giuseppe Velasquez 与 Giuseppe Piazzi 及 Ceres》，油画，约作于 1803 年} \label{fig_Astero_4}
\end{figure}
1772 年，德国天文学家 Johann Elert Bode（援引 Johann Daniel Titius 的观点）发表了一组数列进程，即如今已被否定的 Titius–Bode 定律。除了火星与木星之间存在无法解释的空隙外，该公式似乎能够预测当时已知行星的轨道\(^\text{[24][25]}\)。Bode 对这一“缺失行星”的存在给出了如下解释：

“这一点尤其来源于一项令人惊异的关系，即已知六颗行星在其距日距离上所遵循的规律。若取土星到太阳的距离为 100，则水星距太阳为 4；金星为 4 + 3 = 7；地球为 4 + 6 = 10；火星为 4 + 12 = 16。此处出现一个在规则进程中的空档：在火星之后应为 4 + 24 = 28 的位置，但迄今未见任何行星。难道可以相信宇宙的创造者故意让这一区域保持空缺吗？显然不能。从此处继续，木星距离为 4 + 48 = 52，最终土星为 4 + 96 = 100。”\(^\text{[26]}\)

依据该公式，Bode 预测将在距离太阳约 2.8 天文单位（AU），即约 4.2 亿千米的位置发现另一颗行星\(^\text{[25]}\)。随着 William Herschel 在接近公式预测位置发现了土星外的新行星天王星，Titius–Bode 定律一度得到了强有力支持\(^\text{[24]}\)。1800 年，由德国天文学刊物《Monatliche Correspondenz》（《每月通信》）主编 Franz Xaver von Zach 牵头，他向 24 名资深天文学家发出邀请（并将其称为“天国警察”）\(^\text{[25]}\)，希望他们协同开展针对该“预期行星”的系统性搜寻\(^\text{[25]}\)。虽然他们未能发现 Ceres，但最终陆续发现了编号为 2、3、4 的三颗小行星：2 Pallas、3 Juno 与 4 Vesta\(^\text{[25]}\)。

被选入搜寻小组的天文学家之一是 Giuseppe Piazzi，他是西西里巴勒莫学院的一位天主教神父。在收到加入该小组的正式邀请之前，Piazzi 已于 1801 年 1 月 1 日发现了 Ceres\(^\text{[27]}\)。当时他正在寻找“拉卡伊先生（Mr la Caille）《黄道星目录》中第 87 颗星”\(^\text{[24]}\)，却注意到“其前方又出现了一颗天体”\(^\text{[24]}\)。Piazzi 实际上发现了一颗会移动的、类星状的天体，并最初认为其是一颗彗星\(^\text{[28]}\)：

“其光度略显微弱，颜色类似木星，与通常被视为八等星的许多恒星非常相近。因此，我毫不怀疑其应为一颗普通恒星。[……]然而在第三天夜晚，我的怀疑已转为确信——它绝非一颗固定恒星。不过在公开之前，我仍决定等待至第四天傍晚，当时我欣喜地看到它继续以前几日相同的速率移动。”\(^\text{[24]}\)

Piazzi 共计观测 Ceres 24 次，最后一次记录于 1801 年 2 月 11 日，之后因病中断研究。他于 1801 年 1 月 24 日向两位同行天文学家致信宣布此发现：米兰的同乡 Barnaba Oriani 与柏林的 Bode\(^\text{[21]}\)。在信中，他将其报告为一颗彗星，但同时表示：“鉴于其运动如此缓慢且相当规则，我已多次想到它可能比彗星更为特殊。”\(^\text{[24]}\)

4 月，Piazzi 将完整观测数据寄给 Oriani、Bode 与法国天文学家 Jérôme Lalande，并于当年 9 月发表于《Monatliche Correspondenz》期刊\(^\text{[28]}\)。

其时，由于地球绕太阳公转引起的视位置变化，Ceres 的位置已发生偏移，并过于靠近太阳炫光，其他天文学家无法立即证实 Piazzi 的观测结果。年末 Ceres 理论上应再次可见，但由于间隔时间长、轨道参数尚不精确，其位置难以预测。为重新寻获 Ceres，时年仅 24 岁的数学家 Carl Friedrich Gauss 开发出一种高效的轨道确定方法\(^\text{[28]}\)。仅数周后，他完成预测并将结果寄给 von Zach。1801 年 12 月 31 日，von Zach 与另一位“天国警察”Heinrich W. M. Olbers 在几乎完全符合预测的位置成功重新发现 Ceres\(^\text{[28]}\)。Ceres 位于距太阳 2.8 AU 的轨道位置，与 Titius–Bode 定律几乎完全吻合；然而海王星在 1846 年被发现后，其距离预测值偏差高达 8 AU，促使大多数天文学家最终将该定律视为巧合\(^\text{[29]}\)。Piazzi 将新发现的天体命名为 Ceres Ferdinandea，以“纪念西西里守护女神与波旁王朝国王 Ferdinand”\(^\text{[26]}\)。
\subsubsection{进一步搜寻}
\begin{figure}[ht]
\centering
\includegraphics[width=8cm]{./figures/321bbb2080ca89a0.png}
\caption{首十颗被发现的小行星的尺寸，与月球的比较} \label{fig_Astero_5}
\end{figure}
在随后的数年中，von Zach 小组又陆续发现了另外三颗小行星，即 2 Pallas、3 Juno 与 4 Vesta，其中 Vesta 于 1807 年被发现[25]。此后新的小行星发现陷入停滞，直到 1845 年才出现进展。业余天文学家 Karl Ludwig Hencke 自 1830 年起开始搜寻新小行星，十五年后在寻找 Vesta 的过程中意外发现了后来被命名为 5 Astraea 的小行星，这是 38 年以来首次获得的新发现。Carl Friedrich Gauss 受邀为其命名。此事件促使更多天文学家加入研究，到 1851 年底共计发现 15 颗小行星。1868 年，当 James Craig Watson 发现第 100 颗小行星时，法国科学院在一枚纪念章上雕刻了三位当时最成功的小行星发现者 Karl Theodor Robert Luther、John Russell Hind 与 Hermann Goldschmidt 的肖像，以纪念这一成果[30]。

1891 年，Max Wolf 首创使用天体摄影法（astrophotography）来探测小行星，长时间曝光的照相底片上，小行星会呈现为短条状轨迹[30]。与以往的目视观测相比，这一技术显著提高了发现效率：Wolf 自 323 Brucia 起共发现 248 颗小行星[31]，而此前历史上总共也不过发现略多于 300 颗。尽管天文学界已知小行星数量远不止于此，但多数天文学家对其兴趣不大，甚至有人称其为“天空的害虫”（vermin of the skies）[32]，此说法被分别归因于 Eduard Suess[33] 与 Edmund Weiss[34]。即便在一个世纪后的 20 世纪，仅有数千颗小行星获得识别、编号与命名。
\subsubsection{19 与 20 世纪}
\begin{figure}[ht]
\centering
\includegraphics[width=8cm]{./figures/dc15219005826826.png}
\caption{此处以雷达成像展示的 2013 EC，仍仅具有临时编号} \label{fig_Astero_6}
\end{figure}
过去，小行星的发现通常遵循四个步骤。首先，使用广视场望远镜或天体照相仪对天空某一区域进行拍摄，同一区域需拍摄成对照片，两张照片的拍摄时间间隔通常为一小时，多天内可进行多次成对拍摄。其次，将同一区域的两张底片或胶片置于体视镜（stereoscope）下进行比对，由于绕日运行天体会在两次拍摄期间产生微小位移，其影像在体视镜观察下会呈现为略高于恒星背景的悬浮效果。第三，在确认目标为移动天体后，需使用数字化显微测量仪对其位置进行精确测量，并相对于已知恒星的坐标进行标定[35]。

这前三个步骤尚不能构成对小行星的正式发现：此时观测者仅获得了一次“视象”（apparition），并会被赋予一个临时编号（provisional designation），其格式由发现年份、发现所在半月区间的字母以及按顺序分配的字母与数字组成（示例：1998 FJ74）。最后一步是将观测位置与观测时间提交至小行星中心（Minor Planet Center），由其计算机程序判断该视象是否与以往记录关联并构成同一轨道。若条件满足，该天体将获得正式星表编号，而首次具备可计算轨道视象的观测者将被确认为发现者，并在国际天文学联合会批准下获得命名权[36]。
\subsection{命名}
到 1851 年为止，英国皇家天文学会认为小行星的发现速度过快，需要采用新的分类或命名体系。1852 年，当 de Gasparis 发现第 20 颗小行星时，Benjamin Valz 为其赋予名称和序号，即 20 Massalia，以示其在小行星发现序列中的位置。有些小行星在发现后未能再次被观测到，因此自 1892 年起，新发现的小行星改以发现年份加大写字母的方式记录，用以标识当年同类天体的登记顺序。例如，1892 年发现的前两颗小行星被标记为 1892A 与 1892B。然而，到 1893 年时，字母数已不足以标识全部发现天体，于是在 1893Z 之后使用 1893AA。此后又尝试过多种变体，包括 1914 年起采用“年份 + 希腊字母”的方式。最终，于 1925 年确立了简明的按时间顺序编号体系[30][37]。

目前，所有新发现的小行星均会先获得临时编号（provisional designation），例如 2002 AT4。其格式由发现年份和表示发现半月周期及序号的字母数字组合构成。当小行星轨道获得确认后，将赋予其正式编号，并可在之后获得名称（例如 433 Eros）。正式命名规约使用括号包围编号，如 (433) Eros，但省略括号的写法也相当普遍。在非正式语境中，人们常省略编号；或在连续文本中仅于首次出现时标注编号[38]。此外，小行星的命名可以由发现者按照国际天文学联合会制定的规范提出[39]。
\subsubsection{符号}  
最早被发现的小行星被赋予类似行星符号的标志性象征符号。到 1852 年时，已存在约二十余种小行星符号，且常出现多个变体[40]。

1851 年，在第 15 颗小行星 Eunomia 被发现之后，Johann Franz Encke 在将出版的 1854 年版《柏林天文年鉴》（Berliner Astronomisches Jahrbuch, BAJ）中进行了重大修改。他引入了传统用于表示恒星的圆盘符号（即圆圈）作为小行星的通用符号，并以发现顺序在圆圈中标注数字，以表示特定小行星。该编号圆圈体系迅速被天文学界采用，其后的第 16 颗小行星（1852 年发现的 16 Psyche）成为首个在发现时即采用该体系标示的小行星。然而 Psyche 同时仍被赋予传统象征符号，此后若干小行星亦然。20 Massalia 是首个未获象征符号的小行星，而最后一个被赋予象征符号的小行星是 1855 年发现的 37 Fides[e][41]。
\subsection{形成} 
许多小行星是原行星（planetesimals）的破碎残骸，这些天体曾存在于年轻太阳的星云中，但未能成长至行星尺度[42]。普遍认为，小行星带中的原行星最初与太阳星云中其他天体一样经历演化，直到木星增至接近其当前质量时，由其轨道共振所产生的摄动将带中超过 99\% 的原行星弹出。数值模拟结果，以及自转速率与光谱性质的不连续性均表明：直径大于约 120 km（75 mi）的小行星可能在早期阶段经吸积形成，而更小的天体则是木星扰动期间或其后，小行星之间碰撞所产生的碎片[43]。Ceres 与 Vesta 的体量足够大，因此经历了熔融与分化，重金属元素下沉形成核心，而岩质矿物则构成外壳[44]。

根据 Nice 模型，许多柯伊伯带天体被捕获进入外部小行星带，位于距离太阳大于 2.6 AU 的区域。其中大多数随后被木星再次弹出，但残留者可能构成 D 型小行星，并且可能包括 Ceres 本身[45]。
\subsection{太阳系内的分布}
\begin{figure}[ht]
\centering
\includegraphics[width=6cm]{./figures/e8184c45bcc8665a.png}
\caption{内太阳系中小行星族群位置的俯视示意图} \label{fig_Astero_7}
\end{figure}
在太阳系内部轨道上已发现多类动力学性质不同的小行星族群。它们的轨道受太阳系中其他天体的引力扰动以及雅可夫斯基效应（Yarkovsky effect）的影响。主要族群包括：
\subsubsection{小行星带}
\begin{figure}[ht]
\centering
\includegraphics[width=6cm]{./figures/a627147494f18831.png}
\caption{一幅展示内太阳系行星及小行星族群分布的示意图。距日距离按比例绘制，天体尺寸不按比例。} \label{fig_Astero_8}
\end{figure} 
已知多数小行星位于火星与木星轨道之间的小行星带内，其轨道通常具有较低偏心率（即轨道形状并不显著拉长）。据估计，该区域包含直径大于 1 km（0.6 mi）的小行星约 1.1 至 1.9 百万颗[46]，以及更多体积更小的成员。这些小行星可能是原行星盘的残余物，而在太阳系形成早期，由于木星的强大引力扰动，该区域内的原行星无法通过吸积进一步成长为行星。

与大众印象相反，小行星带大部分区域实际上是空旷的。小行星在极大体积空间内分布稀疏，因此若不进行精确导航，要主动抵达某一目标小行星几乎不可能。尽管如此，当前已知小行星数量达数十万颗，而其总数可能达到数百万甚至更多，这取决于所设定的尺寸下限。目前已知超过 200 颗小行星的直径超过 100 km[47]，且红外波段的观测表明，小行星带内直径大于 1 km 的小行星数量介于 700,000 至 1.7 百万之间[48]。已知小行星的绝对星等多数介于 11 至 19 之间，中位数约为 16[49]。

小行星带的总质量估计约为 \(2.39\times10^{21}\,\mathrm{kg}\)，仅相当于月球质量的约 3\%；相比之下，柯伊伯带与散射盘（Scattered Disk）的总质量超过其 100 倍[50]。在小行星带总质量中，四大天体——Ceres、Vesta、Pallas 与 Hygiea——共占约 62\%，其中 Ceres 单独便占约 39\%。


